\hypertarget{appendix-537171d69d27}{}
\section*{Quantization error metrics}
\section{Quantization error metrics}\label{quantization-error-metrics}

\textbf{Curriculum depth:} Developed · \textbf{First seen:} C006 · \textbf{Latest development:} C006

\subsection{Definition and mechanism}\label{definition-and-mechanism}

For N\textgreater0, RMSE is sqrt(sum squared discrepancies/N), MAE is mean absolute discrepancy, and max error is the largest observed magnitude. Promote float operands before subtraction when measuring in double.

\subsection{Prerequisites}\label{prerequisites}

\begin{itemize}
\tightlist
\item
  \hyperref[appendix-3ae74e89b4e6]{Absolute and relative error} --- supplies the representation or validity rule used above.
\item
  \href{}{Quantization and compressed weights (04-Visual-Computing-II, appendix-fdc7879e8bf2)} --- supplies the representation or validity rule used above.
\end{itemize}

\subsection{Used by}\label{used-by}

Connect later dependent concepts here when their explanations use this mechanism.

\subsection{Recall and next depth}\label{recall-and-next-depth}

Reconstruct the example and identify the condition that would invalidate it.

\textbf{Pending:} Representation error alone does not establish downstream task correctness.

\subsection{Source and navigation}\label{source-and-navigation}

\href{Continuous\%20Notes/Computing\%20Foundations\%20-\%20From\%20Bits\%20to\%20C++\%20and\%20Systems\%20APIs\#Choosing\%20what\%20information\%20to\%20keep}{Continuous teaching story} · \hyperref[appendix-fd090d982a15]{Source coverage ledger} · \href{}{Dependency Map (02-Foundations-II, appendix-4ad47bdd269a)}

\subsection{A metric must match what the consumer loses}\label{a-metric-must-match-what-the-consumer-loses}

Quantization error is the difference \(e=x-\hat x\) between source and reconstruction, but one scalar summary cannot answer every engineering question. Mean squared error emphasizes large deviations; mean absolute error is less sensitive to outliers; maximum absolute error exposes worst-case bounds; relative error becomes useful when scale varies but is unstable near zero.

For tensors, per-tensor MSE can hide one badly damaged channel or group. Diagnostics should often include per-group distributions, saturation counts, scale histograms, and downstream task error. A quantizer that minimizes weight MSE is not guaranteed to minimize model-output error because layers amplify directions differently.

Signal metrics also depend on representation. Comparing gamma-encoded RGB with ordinary squared error is not the same as comparing linear-light energy. Normals may need angular error rather than component MSE.

A useful quantization benchmark therefore reports storage cost together with several reconstruction metrics and one consumer-level metric. The metric is part of the experiment hypothesis, not a decorative number added after compression.
